\hypersetup{pdftitle={Ten Problems That Teach Calculus: Student Pilot Handouts}}
\hero{STUDENT INVESTIGATION \quad 06A \quad / \quad PREDICT}{How far did the cart go?}{Two clients ask about the same trip. Do they need the same number?}
\namefields

\panel{\textbf{The job.} A cart starts at zero on a straight track, rolls backward, then rolls forward past zero. The track technician wants to know where to look for it when the run ends. The wheel-testing technician wants to know how many meters its wheels rolled along the track. Both ask, \emph{How far did it go?}}

\sub{1. Re-create the motion before doing the arithmetic}
Slide a token on a strip of paper. Mark the start, a point to its left where it turns, and a finish to the right of the start. Sketch your trip below. Label the start, turn, and finish and draw arrows for both legs. Your sketch need not be to scale.
\begin{center}
\begin{tikzpicture}[x=1cm,y=1cm]
 \draw[ax] (0,0)--(15.3,0) node[lab,right] {Position};
 \draw[Muted] (7.35,-.15)--(7.35,.15);
 \node[lab,below] at (7.35,-.16) {Start \(0\)};
 \node[lab,above] at (1.18,.65) {Backward};
 \node[lab,above] at (13.75,.65) {Forward};
\end{tikzpicture}
\end{center}

\sub{2. Write a useful reply to each client}
Imagine the cart came back exactly to the zero mark instead. What number would the track technician report? Would the wheel technician report zero too? Say why.
\blank{3}

\sub{3. Decide what to do with a backward reading}
The recorder reports velocity in meters per second. A negative reading means backward. If it reports \(-0.8\uom{m/s}\) throughout a quarter-second interval, what distance and direction does that represent? How should each client count it?
\blank{3}

\panel[Sand]{\textbf{Before you turn the page.} Explain your two counting rules to a partner in ordinary language. Make sure each rule says what to do when the cart goes backward.}
\foot{Student handout \enspace / \enspace 06A}

\clearpage
\hero{STUDENT INVESTIGATION \quad 06B \quad / \quad BUILD AND TEST}{Turn readings into a report}{Make a rule that would still work on somebody else's cart log.}

\begin{center}
\cartplot
\end{center}
{\small\color{Muted}The recorder's velocity model is \(v(t)=t^2-1\) for the first two seconds. Use the representative readings to estimate the contribution of four equal intervals.}

{\setlength{\tabcolsep}{8pt}\renewcommand{\arraystretch}{1.31}
\begin{tabularx}{\linewidth}{@{}l r r >{\raggedright\arraybackslash}X >{\raggedright\arraybackslash}X@{}}\toprule
Interval (s) & Midpoint (s) & Velocity (m/s) & Track report contribution (m) & Wheel report contribution (m)\\\midrule
0--0.5 & 0.25 & \(-0.9375\) & & \\
0.5--1 & 0.75 & \(-0.4375\) & & \\
1--1.5 & 1.25 & \(+0.5625\) & & \\
1.5--2 & 1.75 & \(+2.0625\) & & \\
\midrule
\textbf{Totals} &&&&\\\bottomrule
\end{tabularx}}

\sub{4. Send the report}
State your estimates in sentences with units. Give another team a rule it could use on a new file with unequal interval lengths. Say one way your estimates could improve.
\blank{3}

\sub{5. New job \enspace / \enspace Individual response}
A battery logger reports the currents below. Positive current flows into the battery. A service team needs the net charge gained and the total charge that passed through the meter. What should they report in coulombs? Show how you chose your method. One ampere for one second transfers one coulomb.

{\sffamily\small\color{Muted}Constant current by interval \hfill \(0\)--\(2\) s: \(+3\) A \hfill \(2\)--\(5\) s: \(-1\) A \hfill \(5\)--\(6\) s: \(+2\) A\par}
\blank{3}
\foot{Student handout \enspace / \enspace 06B}

\clearpage
\hero{STUDENT INVESTIGATION \quad 09A \quad / \quad PREDICT AND COMMIT}{When will the drink be ready?}{Give the kitchen crew a prediction it can use before the drink is ready.}
\namefields

\panel{\textbf{The job.} A kitchen crew needs to schedule drinks to reach \(40\tc\). For this scheduling exercise, a predicted time is useful if it is within \textbf{2 minutes} of the time estimated from later readings. Find a procedure another crew can reuse. Your teacher will give you a cup and thermometer or a practice record.}

\sub{1. Predict before you fit a rule}
Sketch your expected drink temperature against time. Show room temperature and \(40\tc\) if your drink starts above \(40\tc\) and the room is cooler than that. Mark when you expect the fastest cooling and explain why.
\begin{center}
\begin{tikzpicture}[x=1cm,y=1cm]
 \draw[step=.5cm,Line!64!white,line width=.35pt] (0,0) grid (14,2.0);
 \draw[ax] (0,0)--(14.3,0) node[lab,right] {Time};
 \draw[ax] (0,0)--(0,2.24) node[lab,right] {Temperature};
\end{tikzpicture}
\end{center}
\blank{2}

\sub{2. Record what you are allowed to use}
Live run: record times 0, 4, and 8 minutes. Use the first and last readings to develop your first rule. Use the 4-minute reading to check its early behavior. Your final check comes from readings you haven't seen. For a practice run, use the values below.

{\setlength{\tabcolsep}{7pt}\renewcommand{\arraystretch}{1.25}
\begin{tabularx}{\linewidth}{@{}l X X X X@{}}\toprule
 & Room & Drink at 0 min & Drink at 4 min & Drink at 8 min\\\midrule
Temperature \(({}^\circ\mathrm C)\) & \rule{.55in}{.35pt} & \rule{.55in}{.35pt} & \rule{.55in}{.35pt} & \rule{.55in}{.35pt}\\
Practice only & \(22.0\) & \(70.0\) & \(61.0\) & \(54.0\)\\\bottomrule
\end{tabularx}}

\sub{3. Commit before seeing the later data}
Write your rule, the predicted crossing time measured from the start, and one assumption. What might the later readings show that would make you change your rule?
\blank{3}
\foot{Student handout \enspace / \enspace 09A \enspace / \enspace Timestamp the prediction}

\clearpage
\hero{STUDENT INVESTIGATION \quad 09B \quad / \quad CHECK AND REVISE}{Will your rule hold up?}{Test it on new readings. Then write a procedure someone else can use.}

\sub{4. Test with later data}
Record the readings your teacher releases. If two consecutive readings straddle \(40\tc\), give their times as a bracket. Estimate the crossing by assuming a straight line between the two readings. If there is no crossing during observation, report what the readings do and do not show.
{\setlength{\tabcolsep}{8pt}\renewcommand{\arraystretch}{1.29}
\begin{tabularx}{\linewidth}{@{}l*{4}{>{\centering\arraybackslash}X}@{}}\toprule
Time from start (min) & \(12\) & \(16\) & \(20\) & \(24\)\\\midrule
Later temperature \(({}^\circ\mathrm C)\) &&&&\\
Rule's predicted temperature \(({}^\circ\mathrm C)\) &&&&\\\bottomrule
\end{tabularx}}
{\small\color{Muted}For a live run, add rows on the back if the drink needs longer to cross.}

\sub{5. Give the client a decision}
State your prediction, the observed bracket, and the estimate of the crossing. Is the estimate within 2 minutes of the prediction? Would every time in the bracket give the same verdict? If not, what new measurement would settle it? Keep your original prediction visible even if you revise the rule.
\blank{4}

\sub{6. Hand off your model}
Write a short procedure another crew can use on a different cup. Include the readings it needs, how it predicts a crossing, and when to distrust the prediction. What should the crew do if the room is already at or above \(40\tc\)?
\blank{3}

\sub{7. New job \enspace / \enspace Individual response}
A room sensor measures \(60\%\) humidity. The long-run level is expected to be \(40\%\). At 5 minutes the reading is \(55\%\). A client asks when it will reach \(45\%\). Use these numbers to make and justify a prediction. What later reading would you request to test your rule?
\blank{3}
\foot{Student handout \enspace / \enspace 09B}
